\subsection{Quasicoherent sheaves}

\subsubsection{The category of quasicoherent sheaves}

Recall that for a ringed space $\O_X$ we have an Abelian category $\Mod_{\O_X}$ of $\O_X$-modules.
Unfortunately this category is often unwieldy, especially when $X$ is a scheme as we would like our modules to also locally look nice.
Recall the embedding $\Mod_A\inj\Mod_{\spec A}$ which maps $M\mapsto\tilde M$, whose sections on the distinguished open sets are
$$ \tilde M(\dof f) = M_f . $$

\bdefn

    If $X$ is a scheme, then an $\O_X$-module $\F$ is called a {\emphcolor quasicoherent sheaf} if for every affine $\spec A\subseteq X$,
    $\F|_{\spec A}\cong\tilde M$ for some $A$-module.
    The full subcategory of $\Mod_{\O_X}$ consisting of quasicoherent sheaves is denoted $\QCoh_X$.

\edefn

\bthrm

    To check that an $\O_X$-module $\F$ is quasicoherent, it is sufficient to check that $\F$ restricts to an $A$-module on any affine cover of $X$.

\ethrm

\bproof

    We will show that the property $\F|_{\spec A}\cong\tilde M$ satisfies the conditions of \refmath{affcomm}.
    The first condition is easier, as if $\F|_{\spec A}\cong\tilde M$, then $\F|_{\spec A_f}\cong\tilde M_f$.

    Now suppose $A=(f_1,\dots,f_n)$, $\F$ is an $\O_{\spec A}$-module such that $\F|_{\dof f_i}\cong\tilde M_i$ where $M_i$ are $A_{f_i}$-modules.
    We want to show that there is an $A$-module $M$ such that $\F\cong\tilde M$.

    Now consider the isomorphisms
    \diagram{
        &\F|_{\dof f_i\cap\dof f_j}\cr
        \tilde M_i|_{\dof f_i\cap\dof f_j}&&\tilde M_j|_{\dof f_i\cap\dof f_j}
    }{
        \diagarrow{from={1,2}, to={2,1}}
        \diagarrow{from={1,2}, to={2,3}}
    }

    By taking sections over $\dof f_i\cap\dof f_j$, this yields an isomorphism $\phi_{ij}\colon(M_i)_{f_j}\to(M_j)_{f_i}$ of $A_{f_if_j}$-modules.
    Furthermore these $\phi_{ij}$ satisfy the cocycle conditions, as is easily verified.
    Let $M_{ij}=(M_i)_{f_j}$ which we identify with $M_{ji}=(M_j)_{f_i}$ via $\phi_{ij}$.

    Now define $M$ to be such that the following sequence is exact (i.e.\ $M=\ker\gamma$):
    $$ 0 \to M \to \prod_{i=1}^n M_i \xtos\gamma \prod_{i<j}M_{ij},\qquad \gamma(m_i)_i = (m_i-m_j)_{i<j} . $$

    Now we claim that $\tilde M_{\dof f_i}\cong\tilde M_i$, i.e.\ $M_{f_i}\cong M_i$.
    Furthermore we claim that these isomorphisms fit into a commutative triangle

    \diagram{
        &M_{f_if_j}\cr
        M_{ij}&M_{ji}
    }{
        \diagarrow{from={1,2}, to={2,1}}
        \diagarrow{from={1,2}, to={2,3}}
        \diagarrow{from={2,1}, to={2,3}, text={\phi_{ij}}, y distance=.25cm}
    }

    Let us assume that $i=1$, then because localization is exact, we have that $M_{f_1}=\ker\gamma_{f_1}$.
    On the other hand the following sequence
    $$ 0 \to M_1 \xtos\beta \prod_{i=1}^n(M_i)_{f_1} \xtos{\gamma_{f_1}} \prod_{i<j}(M_{ij})_{f_1},\qquad \beta(m)=(m_i)_i $$
    is exact.
    This is because we can rewrite it as
    $$ 0 \to M_1 \to \prod_{i=1}^n(M_1)_{f_i} \to \prod_{i<j}(M_i)_{f_if_j} , $$
    which is exact.
    This gives us an isomorphism $M_1\cong M_{f_1}$ as desired.

    Now it is straightforward to prove that these isomorphisms commute with $\phi_{ij}$.
    Thus $\tilde M$ is isomorphic to the gluing of the $\tilde M_i$, but this gluing is $\F$, as desired.
    \qed

\eproof

\bexam

    Let $X=\spec k[t]$, and let $\F$ be the skyscraper sheaf supported at $[(t)]$ with group $k(t)$ and the usual $k[t]$-module structure.
    Because $k(t)$ is the field of fractions of $k[t]$, it has a canonical $k[t]_f$-module structure for all $f\in k[t]$.
    Thus $\F$ is an $\O_X$-module.

    But it is not quasicoherent: if it were then $\F$ would be isomorphic to $\tilde{\Gamma(X,\F)}=\tilde{k(t)}$.
    But this is not the case as $\F(\dof t)=0$ (since $(t)\notin\dof t$), while $k(t)_t\cong k(t)\neq0$.

    In general if $X$ is an integral scheme with function field $K(X)$, then the skyscraper sheaf $p_*K(X)$ is a non-quasicoherent $\O_X$-scheme unless $p=\eta$ is generic.
    We can restrict this to an affine subscheme containing $p$.
    Then the above argument works as well, as there is an $f\in A$ such that $p\notin\dof f$ and so $p_*K(X)(\dof f)=0$ while $\tilde{K(X)}(\dof f)=K(X)_f\cong K(X)\neq0$.
    Meanwhile if $p=\eta$ is generic then $\eta_*K(X)(\dof f)=K(X)$ for all $f\in A$ because $\eta$ lies in all of the open subsets of $X$.

\eexam

Recall that an $A$-module $M$ is {\emphcolor torsion-free} if $am=0$ implies that either $a$ is a zerodivisor or $m=0$.
If $A$ is an integral domain this is equivalent to $a=0$ or $m=0$.
An element $m\in M$ for which there exists a non-zerodivisor $a\in A$ with $am=0$ are called {\emphcolor torsion elements}.
If $M$ is an $A$-module, define $M_\tor$ to be the set of all torsion elements in $M$.
This is a submodule of $M$: if $am,bn=0$ for non-zerodivisors $a,b\in A$ then $ab\in A$ is also a non-zerodivisor and $(ab)(m+n)=0$ so that $m+n\in M_\tor$.
Clearly $M_\tor$ is also closed under scalar multiplication.

\bnote

    In the case of $A$ being an integral domain, the definition of the torsion submodule becomes
    $$ M_\tor = \set{m\in M}[am=0\hbox{ for some nonzero $a\in A$}] . $$
    Sometimes this is taken as the general definition of the torsion subset, and this is not in general a submodule of $M$ as it is not necessarily closed under addition.

    For example let $A=\bZ/\bZ_6$.
    Then $2\cdot3=0$ so they are considered torsion (in this new definition).
    But $2+3=5=-1$, which is not torsion.

    We will basically only care about torsion when $A$ is an integral domain, so we don't need to worry.

\enote

A module is {\emphcolor torsion} if all of its elements are torsion, i.e.\ $M=M_\tor$.
Over an integral domain this is equivalent to $M\otimes_AK(A)=0$.
Indeed if $M$ is torsion, let $m\in M$ then there is a nonzero $a\in A$ such that $am=0$.
Then for all nonzero $b\in A$,
$$ \frac mb = \frac{am}{ab} = 0 $$
so that $M=0$ (recall that $M\otimes_AK(A)\cong M_{(0)}$).
Conversely if $m\in M$ is torsion-free then $m/1\in M\otimes_AK(A)$ is nonzero, as it is zero only when there is a nonzero $a\in A$ for which $am=0$.

\bdefn

    If $X$ is a scheme, an $\O_X$-module is {\emphcolor torsion-free} if $\F_p$ is a torsion-free $\O_{X,p}$-module for all $p\in X$.

    If $X$ is reduced, then a quasicoherent sheaf $\F$ is {\emphcolor torsion} if its stalk at the generic point of every irreducible component is zero.

\edefn

If $X$ is integral then this is the same as saying $\F_p$ is a torsion $\O_{X,p}$-module for all $p\in X$.
Indeed if this is true then $\F_\eta$ is a torsion $\O_{X,\eta}=K(X)$-module, and since $K(X)$ is a field its only torsion modules is the zero module.
Conversely $\O_{X,\eta}$ is the field of fractions of $\O_{X,p}$ for all $p\in X$, and so $\F_\eta$ is the localization of $\F_p$, i.e.\ $\F_\eta=\F_p\otimes K(X)$.
If this is zero then $\F_p$ is torsion.

\bdefn

    If $X$ is a scheme, then its {\emphcolor distinguished affine base} is the subcategory of $\open(X)$ whose objects are the distinguished open subsets of its open affine subsets.
    Then for any $\spec A\subseteq X$ we add only the morphisms $\dof_A(f)\subseteq\dof_A(g)$.

    In other words this category contains all of $X$'s affine open subsets, and for an inclusion $\spec B\subseteq\spec B'$, we add it as a morphism only if $\spec B=\spec A_f$ and $\spec B'=\spec A_g$
    for some $\spec A\subseteq X$.

    Let us denote this category of $\open_\aff(X)$.

\edefn

Note that $\open_\aff(X)^\op$ is a cofinal subposet of $\open(X)^\op$.
If for $p\in X$ we denote by $\open_p(X)\subseteq\open(X)$ the set of neighborhoods of $p$, then $\open_{p,\aff}(X)^\op$ is cofinal in $\open(X)^\op$.

\bremark

    Note here that $\open_\aff(X)$ is not strictly a subposet of $\open(X)$, as it does not include all of its arrows.
    Nevertheless, we can still formulate what it means for it to be cofinal.

    If $\C,\D$ are two categories, a functor $F\colon\C\to\D$ is {\emphcolor final} if for all $d\in D$, the comma category $d/F$ is nonempty and {\emphcolor connected}: all objects can be
    connected by a zig-zag of morphisms.
    It can be shown that this is equivalent to precomposition with $F$ not changing colimits (i.e.\ the canonical $\colim G\circ F\to\colim G$ is an isomorphism).
    For posets the term used is {\emphcolor cofinal} as opposed to final, for historic reasons (the dual of a final functor is {\emphcolor initial}).

    If $\C,\D$ are posets then $d/F$ is the subcategory of $\C$ consisting of elements $c\in C$ such that $d\leq Fc$.
    If $\C$ is filtered then if it is non-empty it is immediately connected.
    And being nonempty just means that $F(\C)\subseteq\D$ is cofinal in the usual sense (every $d\in\D$ is bound by some $c\in\C$).

    In our case $\open_\aff(X)$ and $\open_{p,\aff}(X)$ are both cofiltered.
    If $p\in\spec A\cap\spec B$ then by \refmath[proposition]{affineinter}, $p$ is contained in a common affine neighborhood of $\spec A,\spec B$.
    And clearly $\open_{p,\aff}(X)\subseteq\open_p(X)$ is cofinal, as the affine neighborhoods of $X$ form a base.

\eremark

Because cofinal subposets don't affect colimits, for any sheaf $\F$ on $X$, its stalks can be computed relative to the distinguished affine base:
$$ \F_p \coloneqq \colim\F|_{\open_p(X)^\op} \cong \colim\F|_{\open_{p,\aff}(X)^\op} . $$

Any sheaf on $X$ defines a sheaf on the distinguished affine base.
But conversely any sheaf on $\open_\aff(X)^\op$ defines one on $X$.
The construction is by taking the sheaf $\F^b$ and defining the sheaf $\F$ on an open $U\subseteq X$ by taking the collection of compatible germs:
$$ \F(U) = \set{(f_p\in\F^b_p)_{p\in U}}[\mltext{for all $p\in U$ there is a $U_p\in\open_{p,\aff}(X)$ and a $s\in\F^b(U_p)$\cr such that $s_q=f_q$ for all $q\in U_p$}] . $$
This defines an equivalence between sheaves on $X$ and those defined on the distinguished affine base.

Suppose $\F$ is an $\O_X$-module, then for any distinguished $\spec A_f\inj\spec A$ we have the restriction map $\res\colon\Gamma(\spec A,\F)\to\Gamma(\spec A_f,\F)$.
Because the domain is an $A$-module and the codomain an $A_f$-module, by the universal property of localization this factors through $\Gamma(\spec A,\F)_f$:

\diagram{
    \Gamma(\spec A,\F)&&\Gamma(\spec A_f,\F)\cr
    &\Gamma(\spec A,\F)_f
}{
    \diagarrow{from={1,1}, to={1,3}, text=\res, y distance=.25cm}
    \diagarrow{from={1,1}, to={2,2}}
    \diagarrow{from={2,2}, to={1,3}, text=\alpha, x distance=.25cm}
}

Now $\F$ is quasicoherent iff this $\alpha$ is an isomorphism.
Indeed if $\F$ is quasicoherent, say $\F|_{\spec A}\cong\tilde M$, then this diagram fits into a larger diagram

\diagram{
    \Gamma(\spec A,\F)&&\Gamma(\spec A_f,\F)\cr
    &\Gamma(\spec A,\F)_f\cr
    M&&M_f
}{
    \diagarrow{from={1,1}, to={1,3}, text=\res, y distance=.25cm}
    \diagarrow{from={1,1}, to={2,2}}
    \diagarrow{from={2,2}, to={1,3}, text=\alpha, x distance=.25cm}
    \diagarrow{from={1,1}, to={3,1}, text=\sim, x distance=-.25cm}
    \diagarrow{from={1,3}, to={3,3}, text=\sim, x distance=.25cm}
    \diagarrow{from={3,1}, to={3,3}}
    \diagarrow{from={3,1}, to={2,2}}
    \diagarrow{from={2,2}, to={3,3}, text=\sim, x distance=.25cm}
}

Showing that $\alpha$ is an isomorphism.
On the other hand if $\alpha$ is an isomorphism then it defines an isomorphism $\Gamma(\spec A_f,\F)\to M_f$ making the rightmost triangle commute.
Thus the whole diagram above commutes, showing we have an isomorphism $\F\to\tilde M$.
We summarize this in the following lemma.

\blemm

    An $\O_X$-module $\F$ is quasicoherent iff for every distinguished $\spec A_f\inj\spec A\subseteq X$ the canonical map $\alpha$ below is an isomorphism:

    \diagram{
        \Gamma(\spec A,\F)&&\Gamma(\spec A_f,\F)\cr
        &\Gamma(\spec A,\F)_f
    }{
        \diagarrow{from={1,1}, to={1,3}, text=\res, y distance=.25cm}
        \diagarrow{from={1,1}, to={2,2}}
        \diagarrow{from={2,2}, to={1,3}, text=\alpha, x distance=.25cm}
    }

\elemm

\bdefn

    Let $(X,\O_X)$ be a ringed space.
    Then we say that $f\in\O_X(U)$ is {\emphcolor locally nilpotent} if for every point $p\in U$ there is a neighborhood $p\in V\subseteq U$ such that $f|_V$ is nilpotent
    in $\O_X(V)$.

    We define $X$'s {\emphcolor sheaf of nilpotents} to be the sheaf $\nil(X)$ whose sections are given by
    $$ \O_{\nil X}(U) = \set{f\in\O_X(U)}[\hbox{$f$ is locally nilpotent}] . $$

\edefn

Note that $f\in\O_X(U)$ is locally nilpotent iff for every $p\in U$, $f_p$ is in $\nil(\O_{X,p})$.
This is because $f\in\nil(\O_{X,p})$ iff there is a neighborhood $p\in V$ such that $f|_V\in\nil(\O_X(V))$.
This clearly shows that $\nil(X)$ is an $\O_X$-module, as $\nil(\O_{X,p})$ is an ideal for all $p$.

\bnote

    $\nil(X)$ is an example of an {\emphcolor ideal sheaf} (of $X$).
    This is a subobject of $X$ in $\Mod_{\O_X}$.
    Because $\Mod_{\O_X}$ is complete and its limits are taken pointwise, by an argument identical to that in \refmath[remark]{monepiccomp} a morphism of $\O_X$-modules
    is monic iff each component is.
    Thus an ideal sheaf of $X$ is just a subsheaf $I$ whose sections are all ideals of $X$'s sections: $\O_I(U)\leq\O_X(U)$.

\enote

Recall that the tensor product of $\O_X$-modules is defined as the sheafification of the presheaf whose sections are $\F(U)\otimes_{\O_X(U)}\G(U)$.
We will show that when $\F,\G$ are quasicoherent sheaves on a scheme $X$, there is no need to sheafify.
That is, the presheaf tensor is already a sheaf.

For an affine $\spec A\subseteq X$ suppose $\Gamma(\spec A,\F)=M$ and $\Gamma(\spec A,\G)=N$.
Then define $\Gamma(\spec A,\F\otimes\G)=M\otimes_AN$.
And if $\spec A_f\inj\spec A$ is distinguished then the restriction map $\Gamma(\spec A,\F\otimes\G)\to\Gamma(\spec A_f,\F\otimes\G)$ corresponds to the localization map
$M\otimes_A N\to M_f\otimes_{A_f} N_f\cong(M\otimes_AN)_f$.
Thus $(\F\otimes\G)|_{\spec A}\cong\tilde{M\otimes_AN}$, so that $\F\otimes\G$ is a sheaf on all the open affine subsets of $X$.
In particular $\F\otimes\G$ forms a sheaf on the distinguished affine base of $X$, and thus is itself a sheaf on $X$.

\bnote

    Here we used the fact that $M_f\otimes_{A_f}N_f\cong(M\otimes_AN)_f$.
    In general we have
    $$ S^{-1}M\otimes_{S^{-1}A}S^{-1}N \cong M\otimes_AS^{-1}A\otimes_{S^{-1}A}S^{-1}N \cong M\otimes_AN\otimes_AS^{-1}A \cong S^{-1}(M\otimes_AN) . $$

\enote

\bdefn

    If $X$ is a scheme and $f\in\O_X(X)$ is a global section, let $X_f\subseteq X$ be the subset on which $f$ doesn't vanish.
    This is an open subset and as such an open subscheme of $X$.

\edefn

Note that for $f\in A$, $(\spec A)_f$ and $\spec A_f=\dof f$ are equal.
Indeed, $f$ doesn't vanish at $\p$ iff $f_\p\notin\m_\p=\p A_\p$, i.e.\ iff $f\notin\p$.

\blemm[title={The Qcqs lemma}, name={qcqslemm}]

    Suppose $X$ is a quasicompact quasiseparated scheme, $\F$ a quasicoherent sheaf on $X$, and $f\in X$.
    Then the restriction map
    $$ \res_{X_f\subseteq X} \colon \Gamma(X,\F) \to \Gamma(X_f,\F) $$
    is precisely the localization by $f$ map.

\elemm

\bproof

    Because $X$ is qcqs it is covered by finitely many affine subschemes, the pairwise intersections of which are covered by finitely many affine subschemes.
    So suppose $X=\bigcup_iU_i$ is a finite affine covering, and $U_i\cap U_j=\bigcup_kU_{ijk}$ is a finite affine covering for all $i,j$.

    The sheaf condition on $\F$ necessitates that the following sequence is exact (the arrows are the standard arrows):
    $$ 0 \to \Gamma(X,\F) \to \bigoplus_i\Gamma(U_i,\F) \to \bigoplus_{i,j,k}\Gamma(U_{ijk},\F) . $$
    This is because the direct sums are finite and thus are equivalently products.
    Localization by $f$ is exact and commutes with direct sums, so the following sequence is also exact:
    $$ 0 \to \Gamma(X,\F)_f \to \bigoplus_i\Gamma(U_i,\F)_f \to \bigoplus_{i,j,k}\Gamma(U_{ijk},\F)_f . $$

    Now for an affine $\spec A\subseteq X$, $\F|_{\spec A}\cong\tilde M$ for some $A$-module $M$.
    Then
    $$ \Gamma((\spec A)_f,\F) = \Gamma(\spec A_f,\F) = M_f . $$
    Thus we have an exact sequence
    $$ 0 \to \Gamma(X,\F)_f \to \bigoplus_i\Gamma((U_i)_f,\F) \to \bigoplus_{i,j,k}\Gamma((U_{ijk})_f,\F) $$
    and the last arrow is the usual one induced from restrictions.

    But the kernel of this map is $\Gamma(X_f,\F)$ as per the sheaf axiom.
    Thus we get the desired result.
    \qed

\eproof

Our goal now is to show that $\QCoh_X\subseteq\Mod_{\O_X}$ is an Abelian category.
In general if $\B\subseteq\A$ is a full subcategory of an Abelian category $\A$, then for $\B$ to be an Abelian category itself it is sufficient that
\benum
    \item $\B$ contain $\A$'s zero element,
    \item $\B$ be closed under finite sums,
    \item $\B$ be closed under kernels and cokernels.
\eenum
Indeed, if $\B$ is closed under colimits taken in $\A$, then those are $\B$'s colimits as well.
Thus if $f$ is monic in $\B$, then $\ker f=0$ as the kernel of a monic is zero, and thus $f=\ker\coker f$ in $\A$ and thus in $\B$.
Similar for epimorphisms.

Clearly the zero $\O_X$-module is quasicoherent.
And (2) follows from the fact that $\tilde M\oplus\tilde N\cong\tilde{M\oplus N}$, so the sum of quasicoherent sheaves is quasicoherent.
For (3) if $\alpha\colon\F\to\G$ is a morphism of quasicoherent sheaves, then let $U\subseteq X$ be affine, and denote by $\beta\colon M\to N$ the induced morphism of $A$-modules.
We then {\it define\/} $\ker\alpha(U)=\ker\beta$ and $\coker\alpha(U)=\coker\beta$.
Because localization is exact, it commutes with kernels and cokernels, we see that on $U$, $\ker\alpha\cong\tilde{\ker\beta}$ and $\coker\alpha\cong\tilde{\coker\beta}$.
This definition forms a sheaf on the distinguished affine base, so $\ker\alpha,\coker\alpha$ are quasicoherent sheaves.

Recall that exactness can be checked stalk-locally.
Thus it can be checked on an affine cover, which shows that
$$ 0 \to \ker\alpha \to \F \xtos\alpha \G \to \coker\alpha \to 0 $$
is exact, as it resorts to checking on affine subschemes, which just gives
$$ 0 \to \ker\beta \to M \xtos\beta N \to \coker\beta \to 0 . $$

So we have proven the following.

\bthrm

    $\QCoh_X$ is an Abelian subcategory of $\Mod_{\O_X}$.

\ethrm

For the last part of this section, we will show an alternative method of defining quasicoherent sheaves, such that the definition works for any ringed space.
Because manifolds are also examples of ringed spaces, this allows us to define quasicoherent sheaves on a manifold.

\bprop

    Let $A$ be a ring, then there is an adjunction
    $$ \Mod_A \adjarrow{\tilde\bul}{\Gamma} \Mod_{\O_{\spec A}} $$
    between $M\mapsto\tilde M$ and global sections $\F\mapsto\Gamma(\spec A,\F)$.

\eprop

\bproof

    This just amounts to showing that global sections
    $$ \hom_{\spec A}(\tilde M,\F) \xtos\Gamma \hom_A(M,\Gamma\F) $$
    is bijective.
    This is clear because if $\phi\colon M\to\Gamma\F$ is a module map then for any $f\in A$ by localization there is a unique $\phi_f$ making the following diagram commute

    \commsq{M&\Gamma(\spec A,\F)\cr M_f&\Gamma(\spec A_f,\F)}
        {\phi}{\phi_f}{}{}

    It is easy to see that the $\phi_f$ assemble into a $\O_{\spec A}$-module morphism.
    Existence and uniqueness shows that $\Gamma$ is a bijection.
    \qed

\eproof

In particular $M\mapsto\tilde M$ commutes with colimits.

\bprop

    Let $X$ be a scheme.
    An $\O_X$-module is quasicoherent iff there exists an open cover $X=\bigcup_iU_i$ such that for all $i$, $\F|_{U_i}$ fits into an exact sequence
    $$ \O_{U_i}^{\oplus I} \to \O_{U_i}^{\oplus J} \to \F|_{U_i} \to 0 . $$

\eprop

\bproof

    If $\F$ is a quasicoherent sheaf take $U_i$ to be an affine open cover.
    For an affine $\spec A\subseteq X$ we have that
    $$ \O_{\spec A}^{\oplus I} \cong (\tilde A)^{\oplus I} \cong \tilde{A^{\oplus I}} . $$
    Now if $\F|_{\spec A}\cong\tilde M$ then $M$ fits into an exact sequence
    $$ A^{\oplus I} \to A^{\oplus J} \to M \to 0 . $$
    Because $\tilde\bul$ is a left-adjoint and thus right-exact, we have that the following sequence
    $$ \matrix{
        \tilde{A^{\oplus I}} &\to& \tilde{A^{\oplus J}} &\to& \tilde M &\to& 0\cr
        \cong && \cong && \cong\cr
        \O_{\spec A}^{\oplus I} &\to& \O_{\spec A}^{\oplus J} &\to& \F|_{\spec A} &\to& 0
    } $$
    is exact, as desired.

    Conversely if such an open cover exists, then we can refine it to be an affine cover.
    Because restricting $\O_X\mapsto\O_U$ is exact, we can just assume the cover to be affine.
    Thus we have for every $\spec A\subseteq X$ in the cover, an exact sequence
    $$ \tilde{A^{\oplus I}} \to \tilde{A^{\oplus J}} \to \F|_{\spec A} \to 0 . $$
    Let $M=\coker(A^{\oplus I}\to A^{\oplus J})$, so we have an exact sequence
    $$ A^{\oplus I} \to A^{\oplus J} \to M \to 0 . $$
    Applying $\tilde\bul$ then gives us the exact sequence
    $$ \tilde{A^{\oplus I}} \to \tilde{A^{\oplus J}} \to \tilde M \to 0 , $$
    and the first arrow is the same as the previous one because $\tilde\bul$ and global sections are inverses from $A$-modules.

    Thus we must have that $\F|_{\spec A}\cong\tilde M$, meaning $\F$ is quasicoherent as desired.
    \qed

\eproof

Thus for a general ringed space $(X,\O_X)$ we can take the condition in the above proposition as the definition of quasicoherence.

\subsubsection{Finiteness of quasicoherent sheaves}

For an $A$-module $M$, there are three finiteness conditions we care about:
\benum
    \item $M$ is {\emphcolor finitely generated}: it has a finite number of generators.
    Equivalently there is a surjection $A^{\oplus n}\to M$ for a finite $n$.
    \item $M$ is {\emphcolor finitely presented}: it has a finite number of generators and a finite number of relations.
    Equivalently there is a {\emphcolor finite presentation}: an exact sequence
    $$ A^{\oplus m} \to A^{\oplus n} \to M \to 0 . $$
    \item $M$ is {\emphcolor coherent}: it is finitely generated, and for any (not necessarily surjective) map $A^{\oplus n}\to M$, its kernel is finitely generated.
\eenum

Note that being finitely presented is equivalent to there being a surjection $A^{\oplus n}\to M$ whose kernel is finitely generated.
Indeed if $A^{\oplus m}\to A^{\oplus n}\to M\to0$ is exact, then the kernel of $A^{\oplus n}\to M$ is the image of a map from $A^{\oplus m}$, which is finitely generated.
Conversely if $A^{\oplus n}\to M$ has finitely generated kernel $K$, take a surjection $A^{\oplus m}\to K$; the composition $A^{\oplus m}\to A^{\oplus n}\to M\to0$ is exact.

Thus coherence implies finitely presented.
Clearly finitely presented implies finitely generated.
We will show that if $A$ is Noetherian, finitely generated implies coherence.

\bprop

    If $A$ is Noetherian, all three finiteness conditions above are equivalent.

\eprop

\bproof

    Suppose $M$ is a finitely generated $A$-module, and take a map $\alpha\colon A^{\oplus m}\to M$.
    Then $\ker\alpha$ is a submodule of $A^{\oplus m}$, which is Noetherian as the finite sum of Noetherian modules, and thus $\ker\alpha$ is Noetherian.
    In particular $\ker\alpha$ is finitely generated.
    \qed

\eproof

If $A$ is Noetherian, then the category of finitely generated $A$-modules is an Abelian subcategory of $\Mod_A$.
Indeed, the zero module is finitely generated, and the sum of two finitely generated $A$-modules is finitely generated.
We just need to show that the kernel and cokernel of a map between finitely generated $A$-modules is finitely generated.
The cokernel is a quotient of its codomain and is thus finitely generated.
The kernel is a submodule of the domain; since the domain is finitely generated, it is Noetherian so its submodules are finitely generated.

In fact even when $A$ is not Noetherian, the category of coherent $A$-modules is an Abelian subcategory of $\Mod_A$.
We ignore the proof.

\bprop

    The category of coherent $A$-modules is a full Abelian subcategory of $\Mod_A$.

\eprop

Note that the localization of finitely generated modules are finitely generated: if $M$ is a finitely generated $A$-module, then $S^{-1}M$ is a finitely generated $S^{-1}A$-module.
Conversely if $A=(f_1,\dots,f_n)$ and $M_{f_i}$ is a finitely generated $A_{f_i}$-module, then $M$ is a finitely generated $A$-module.
Suppose $m_{ij}\in M$ generate $M_{f_i}$, then all the $m_{ij}$ generate $M$.
Indeed, for $m\in M$ there are $a_{ij}\in A$ and $k_i\geq0$ such that for all $i$:
$$ m = \frac{\sum_j a_{ij}m_{ij}}{f_i^{k_i}} \iff f_i^{N_i}\parens{f_i^{k_i}m - \sum_j a_{ij}m_{ij}} $$
for some $N_i\geq0$.
We can assume that all the $N_i$ and $k_i$ are equal, and by updating $a_{ij}$ to be $f_i^{N_i}a_{ij}$ we get that
$$ f_i^Nm = \sum_ja_{ij}m_{ij} . $$
Because $(f_i)_i$ is the unit ideal, so too is $(f_i^N)_i$ (because its radical is $(f_i)_i$ and $\sqrt I=(1)$ iff $I=(1)$).
Thus there are $c_i\in A$ such that $\sum_ic_if_i^N=1$ and thus
$$ m = \sum_ic_if_i^Nm = \sum_{i,j}c_ia_{ij}m_{ij} $$
as desired.

We summarize this as follows.

\bprop

    Let $M$ be an $A$-module.
    \benum
        \item If $M$ is a finitely generated $A$-module, then $S^{-1}M$ is a finitely-generated $S^{-1}A$-module.
        \item If $A=(f_1,\dots,f_n)$ and $M_{f_i}$ is a finitely-generated $A_{f_i}$-module for all $i$, then $M$ is a finitely generated $A$-module.
    \eenum

\eprop

\blemm

    Suppose $M$ is a finitely presented $A$-module.
    Then for any surjection $\phi\colon A^{\oplus p}\to M$, its kernel is finitely generated.

\elemm

\bproof

    Choose a finite presentation of $M$:

    \diagram{
        &&A^{\oplus p}\cr
        A^{\oplus m}&A^{\oplus n}&M&0\cr
        &&0
    }{
        \diagarrow{from={2,1}, to={2,2}, text=\beta, y distance=.25cm}
        \diagarrow{from={2,2}, to={2,3}, text=\alpha, y distance=.25cm}
        \diagarrow{from={1,3}, to={2,3}, text=\phi, x distance=.25cm}
        \diagarrow{from={2,3}, to={2,4}}
        \diagarrow{from={2,3}, to={3,3}}
    }

    Because the rings are free, we can find lifts $g,h$ making the following diagram commute

    \diagram{
        &&A^{\oplus p}\cr
        A^{\oplus m}&A^{\oplus n}&M&0\cr
        &&0
    }{
        \diagarrow{from={2,1}, to={2,2}, text=\beta, y distance=.25cm}
        \diagarrow{from={2,2}, to={2,3}, text=\alpha, y distance=.25cm}
        \diagarrow{from={1,3}, to={2,3}, text=\phi, x distance=.25cm}
        \diagarrow{from={2,3}, to={2,4}}
        \diagarrow{from={2,3}, to={3,3}}
        \diagarrow{from={1,3}, to={2,2}, text=g, x distance=-.25cm, x off=-.1cm}
        \diagarrow{from={2,2}, to={1,3}, text=h, x distance=.25cm, x off=.1cm}
    }

    Explicitly, if $x_1,\dots,x_p$ generate $A^{\oplus p}$ then by surjectivity there are $y_i\in A^{\oplus n}$ such that $\alpha(y_i)=\phi(x_i)$.
    So define $g(x_i)=y_i$.

    Now consider the following sequence
    $$ A^{\oplus m}\oplus A^{\oplus k} \xtos\psi A^{\oplus k} \xtos\phi M \to 0,\qquad \psi = (h\circ\beta)\oplus(\id-h\circ g) . $$
    We claim that this sequence is exact, and then we are finished.
    Firsly,
    $$ \phi\circ h\circ\beta = \alpha\circ\beta = 0,\qquad \phi - \phi\circ h\circ g = \phi - \alpha\circ g = \phi - \phi = 0 $$
    so that $\phi\circ\psi=0$.
    Now if $x\in\ker\phi$ we want to find a $y\in A^{\oplus m}$ such that
    $$ x = \psi(y+x) = h(\beta(y)) + x - h(g(x)) \iff h(\beta(y)) = h(g(x)) . $$
    Now if $\phi x=0$ then $\alpha gx=0$ so $gx\in\ker\alpha=\im\beta$, so there is a $y\in A^{\oplus m}$ such that $\beta y=gx$.
    Thus $h\beta y=hgx$ as desired.
    \qed

\eproof

We already showed that finitely generated modules satisfy the hypothesis of the affine communication lemma.
We now show this for finite presentations and coherence.

\blemm

    Let $M$ be an $A$-module.
    \benum
        \item If $M$ is finitely presented, then $S^{-1}M$ is finitely presented as an $S^{-1}A$-module.
        \item If $(f_1,\dots,f_n)=A$ and $M_{f_i}$ is a finitely presented $A_{f_i}$-module for all $i$, then $M$ is a finitely presented $A$-module.
    \eenum

\elemm

\bproof

    (1) is clear because localization is exact.
    For (2) we know then that $M$ is finitely generated.
    Thus there is a surjection $\phi\colon A^{\oplus n}\to M$, and its localizations $\phi_{f_i}\colon A_{f_i}^{\oplus n}\to M_{f_i}$ are all surjective as well (as localization is exact).
    By the above lemma, this means that $\ker(\phi_{f_i})$ is finitely generated.
    Because $\ker(\phi_{f_i})\cong(\ker\phi)_{f_i}$, this means that $(\ker\phi)_{f_i}$ is finitely generated for all $i$, and thus $\ker\phi$ is.
    \qed

\eproof

For coherence we note that $M$ is coherent iff it is finitely generated and all of its finitely generated submodules are finitely presented.
Indeed if $M$ is coherent and $N\subseteq M$ is a finitely generated submodule, let $A^{\oplus n}\to N$ be a surjection.
Then the composition $A^{\oplus n}\to N\inj M$ has a finitely generated kernel, so $N$ is finitely presented.
Conversely suppose all of $M$'s finitely generated submodules are finitely presented.
Let $\phi\colon A^{\oplus n}\to M$ be any map, then consider its kernel $K$.
Then $A^{\oplus n}/K$ is a submodule of $M$, thus finitely presented, and so by the above lemma the kernel of $A^{\oplus n}\to A^{\oplus n}/K$, which is $K$, is finitely generated.

In particular a finitely generated submodule of a coherent module is coherent.

\blemm

    Let $M$ be an $A$-module.
    \benum
        \item If $M$ is coherent, then $S^{-1}M$ is coherent as an $S^{-1}A$-module.
        \item If $(f_1,\dots,f_n)=A$ and $M_{f_i}$ is a coherent $A_{f_i}$-module for all $i$, then $M$ is a coherent $A$-module.
    \eenum

\elemm

\bproof

    For (1) if $N'\subseteq S^{-1}M$ is a finitely generated submodule, then $N'=S^{-1}N$ for a finitely generated submodule $N\subseteq M$ (if $x_i/s_i$ generate $N'$
    then take $N$ generated by $x_i$).
    Then since $M$ is coherent, $N$ is finitely presented, and so by the previous lemma $N'=S^{-1}N$ is finitely presented as well.
    Thus $S^{-1}M$ is coherent.

    For (2) let $N\subseteq M$ be finitely generated.
    Then $N_{f_i}\subseteq M_{f_i}$ is finitely generated, thus finitely presented.
    This is true for all $i$ so by the previous lemma, $N$ is finitely generated.
    Thus $M$ is coherent.
    \qed

\eproof

\bdefn

    A quasicoherent sheaf $\F$ is {\emphcolor finite type} (resp.\ {\emphcolor finitely presented}, {\emphcolor coherent}) if for every affine open $\spec A\subseteq X$,
    $\Gamma(\spec A,\F)$ is a finitely generated (resp.\ finitely presented, coherent) $A$-module.

\edefn

Note that coherent quasicoherent sheaves are finitely presented, which in turn are finitely generated.
Because these three notions are equivalent over a Noetherian ring, these three notions are equivalent over a locally Noetherian scheme.

Because coherent modules are an Abelian subcategory of $A$-modules, coherent quasicoherent sheaves form an Abelian subcategory of $\QCoh_X$.
This can be seen because our construction of the kernel and cokernel of $\alpha\colon\F\to\G$ is the (co)kernel of $\beta\colon M\to N$, and $M,N$ are coherent so the (co)kernel
of $\beta$ is as well.

Just like how quasicoherence can be generalized to arbitrary ringed spaces, so too can finite type, finitely presented, and coherence.

\bdefn

    Let $(X,\O_X)$ be a ringed space and $\F$ an $\O_X$-module.
    \benum
        \item $\F$ is {\emphcolor finite type} if for every $x\in X$ there is an open $x\in U\subseteq X$ and a surjection
        $$ \O_U^{\oplus n} \to \F|_U \to 0 $$
        for some $n\geq0$.
        \item $\F$ is {\emphcolor finitely presentable} if for every $x\in X$ there is an open $x\in U\subseteq X$ and $n,m\geq0$ and an exact sequence
        $$ \O_U^{\oplus m} \to \O_U^{\oplus n} \to \F|_U . $$
        \item $\F$ is {\emphcolor coherent} if $\F$ is finite type and for all open $U\subseteq X$, the kernel of any map
        $$ \O_U^{\oplus n} \to \F|_U $$
        is of finite type.
    \eenum

\edefn

